\documentclass[11pt,a4paper]{article}
\usepackage[margin=25mm,headheight=14pt]{geometry}
\usepackage{fontspec}
\setmainfont{lmroman10-regular.otf}[BoldFont=lmroman10-bold.otf,ItalicFont=lmroman10-italic.otf,BoldItalicFont=lmroman10-bolditalic.otf]
\setsansfont{lmsans10-regular.otf}
\setmonofont{lmmono10-regular.otf}[Scale=MatchLowercase]
\newfontfamily\authorsymbol{Noto Sans Symbols}[Scale=MatchUppercase]
\usepackage{amsmath,amssymb,amsthm,mathtools,booktabs,array}
\usepackage{microtype}
\usepackage[hidelinks,unicode]{hyperref}
\usepackage{fancyhdr}
\pagestyle{fancy}\fancyhf{}
\fancyhead[L]{\small\itshape Two sparse binary operators}
\fancyhead[R]{\small\thepage}\renewcommand{\headrulewidth}{0pt}
\setlength{\parindent}{1.2em}\setlength{\parskip}{0pt}
\setlength{\emergencystretch}{2em}
\newtheorem{theorem}{Theorem}
\newtheorem{proposition}[theorem]{Proposition}
\newtheorem{lemma}[theorem]{Lemma}
\newtheorem{corollary}[theorem]{Corollary}
\theoremstyle{definition}\newtheorem{definition}[theorem]{Definition}
\newcommand{\F}{\mathbb F_2}
\newcommand{\Mat}{\operatorname{M}}
\DeclareMathOperator{\rank}{rank}
\DeclareMathOperator{\spanlin}{span}
\hypersetup{pdftitle={Two Sparse Binary Operators with Optimal Generating Length},pdfauthor={lmtlssss🜐},pdfsubject={Mathematics: matrix algebras over finite fields},pdfkeywords={matrix algebra, generating length, finite field, rank-one certificate, cyclic vector}}
\title{\LARGE Two Sparse Binary Operators\\with Optimal Generating Length}
\author{lmtlssss{\authorsymbol 🜐}}
\date{27 October 2026\\[3pt]{\small\href{https://doi.org/10.5281/zenodo.23701587}{doi:10.5281/zenodo.23701587}}}
\begin{document}
\maketitle\thispagestyle{plain}
\begin{abstract}
We prove that two specified rotation-and-shift operators generate the full
matrix algebra on a 32-dimensional binary vector space. Their generating
length is ten: words of length at most ten span the algebra, and no pair of
32-by-32 matrices can attain a smaller length. An exact certificate gives
1,024 independent words. A second certificate produces a rank-one matrix
as a polynomial in the generators and gives two cyclic bases. It yields
an explicit expansion of every linear map and proves the same generation
statement over every field of characteristic two. We include the finite
certificates, their independent verification, and the consequences for
common invariant subspaces, commutants, and simultaneous linear quotients.
\end{abstract}
\noindent\textit{Keywords.} Matrix algebras; binary linear maps; generating
length; rank-one certificates; cyclic vectors.

\section{The two operators}
Let $V=\F^{32}$, with coordinates $x_0,\ldots,x_{31}$ numbered from the
least significant bit. Define cyclic right rotation and zero-filled right
shift by
\begin{equation}\label{eq:shifts}
 (R_rx)_i=x_{(i+r)\bmod32},\qquad
 (S_sx)_i=\begin{cases}x_{i+s},&i+s<32,\\0,&i+s\geq32.\end{cases}
\end{equation}
All vector and matrix sums below are over the stated field. Set
\begin{equation}\label{eq:operators}
 A=R_7+R_{18}+S_3,\qquad B=R_{17}+R_{19}+S_{10}.
\end{equation}
Each matrix has at most three nonzero entries in each row. These are fixed,
previously specified word operators~\cite[Section~5.1]{rfc6234}; the question
here is the algebra and word length of this particular pair.

For a field $K$ containing $\F$, regard the same binary matrices as elements
of $\Mat_{32}(K)$. Let $L_d(A,B;K)$ be the $K$-linear span of all products
of at most $d$ factors chosen from $A,B$. The empty product is the identity.
The \emph{generating length} is the least $d$ for which this span equals the
unital algebra generated by $A,B$. This convention is standard in the
matrix-algebra length literature~\cite{laffey,klep}.

\begin{theorem}\label{thm:main}
For every field $K$ of characteristic two,
\begin{equation}\label{eq:main}
 K[A,B]=\Mat_{32}(K),\qquad \ell(A,B;K)=10.
\end{equation}
Ten is the smallest generating length attainable by any pair that generates
$\Mat_{32}(K)$. Over $\F$, every one of the $2^{1024}$ linear maps on $V$
is an XOR sum of at most 1,024 basis words, each of length at most ten.
\end{theorem}
The two parts use distinct certificates. Section~2 proves generation by a
rank-one construction. Section~3 certifies the sharp length. Neither
statement identifies a general linear map with a single ten-factor product.

\section{An explicit rank-one certificate}
The following elementary criterion is a constructive form of the standard
rank-one method for full matrix algebras~\cite{radjabalipour}.

\begin{lemma}\label{lem:cyclic}
Let $\mathcal A\subseteq\Mat_n(K)$ be a unital algebra containing $T$ and
$P=uv^{\mathsf T}$. Suppose
\[
 U=(u,Tu,\ldots,T^{n-1}u),\qquad
 V=\begin{pmatrix}v^{\mathsf T}\\v^{\mathsf T}T\\\vdots\\v^{\mathsf T}T^{n-1}\end{pmatrix}
\]
are invertible. Then the $n^2$ matrices $T^iPT^j$, $0\leq i,j<n$, form
a basis of $\Mat_n(K)$, and $\mathcal A=\Mat_n(K)$.
\end{lemma}
\begin{proof}
For a coefficient matrix $H=(h_{ij})$,
\begin{equation}\label{eq:expansion}
 \sum_{i,j=0}^{n-1} h_{ij}T^iPT^j=UHV.
\end{equation}
The map $H\mapsto UHV$ is invertible. Every summand belongs to
$\mathcal A$, and the $n^2$ coefficient matrices with one nonzero entry
therefore give a basis.
\end{proof}

For the operators in~\eqref{eq:operators}, put
\begin{equation}\label{eq:D}
 D=B^3+A,
\end{equation}
and define
\begin{equation}\label{eq:p}
 p(t)=t^{31}+t^{29}+t^{27}+t^{25}+t^{19}+t^{11}+t^9+t.
\end{equation}
Encode a binary column $z$ by the integer $\sum_{i=0}^{31}z_i2^i$.
The two vectors in the certificate are
\begin{equation}\label{eq:uv}
 u=\texttt{aeb51dbb},\qquad v=\texttt{ddb8ad75}
 \quad\text{(hexadecimal)}.
\end{equation}

\begin{proposition}\label{prop:certificate}
With the preceding coordinate convention,
\begin{equation}\label{eq:rankone}
 p(D)=uv^{\mathsf T},\qquad
 \det(u,Au,\ldots,A^{31}u)=1,\qquad
 \det\begin{pmatrix}v^{\mathsf T}\\v^{\mathsf T}A\\\vdots\\v^{\mathsf T}A^{31}\end{pmatrix}=1.
\end{equation}
These are identities over $\F$ and over every field of characteristic two.
\end{proposition}
\begin{proof}
This is an explicit finite certificate. Construct $A,B$ by
\eqref{eq:shifts}--\eqref{eq:operators}, form $D$, and accumulate the eight
powers in~\eqref{eq:p}. For column $j$ the result is $u$ when $v_j=1$ and
zero when $v_j=0$. This checks all 1,024 entries of the first identity.
For the two cyclic matrices $U,V$, Appendix~A gives the columns of
matrices $X,Y$. Exact multiplication gives
\begin{equation}\label{eq:inverses}
 UX=XU=I,\qquad VY=YV=I.
\end{equation}
These products check the two determinant claims, since the only nonzero
scalar of $\F$ is one. All entries and identities use only zero, one,
addition and multiplication. They therefore remain valid under scalar
extension to any field of characteristic two. The companion verifier
reconstructs these calculations in coordinate-row form independently of
the column implementation.
\end{proof}

Taking $T=A$ in Lemma~\ref{lem:cyclic} proves the generation statement in
Theorem~\ref{thm:main}. It also gives the following concrete synthesis rule:
\begin{equation}\label{eq:synthesis}
 H=U^{-1}MV^{-1},\qquad
 M=\sum_{i,j=0}^{31}h_{ij}A^ip(B^3+A)A^j.
\end{equation}
The coefficients are unique. The rank-one matrix in this construction is
nilpotent: $v^{\mathsf T}u=0$, so $P^2=0$. The proof does not treat $P$ as
a projection. The products in~\eqref{eq:synthesis} are a constructive basis;
they are separate from the shorter word basis used next.

\section{The optimal word length}
A word record stores symbols in input order. Thus the record
\texttt{AB} means that $A$ acts first, and its matrix is $BA$.
For a word $w$, append a symbol $g\in\{A,B\}$ by setting
$M_{wg}=gM_w$. Start with the empty word and $M_{\varnothing}=I$.

The supplied \texttt{short-word-certificate.json} lists 1,024 word records.
Their maximum length is ten. Reconstructing each record by the preceding
rule and flattening its matrix gives 1,024 independent vectors in
$\F^{1024}$. One exact way to select this list is to process independent
words in breadth-first order, append $A$ then $B$, and retain a new word
precisely when its matrix increases the current span. Previously rejected
matrices need not be extended: left multiplication preserves every linear
relation that caused their rejection.

For additional transparency, the verifier enumerates \emph{all} 2,047
words of length at most ten. Their cumulative dimensions are:
\begin{center}
\small
\begin{tabular}{@{}rrrrrrrrrrrr@{}}
\toprule
$d$&0&1&2&3&4&5&6&7&8&9&10\\
\midrule
$\dim L_d$&1&3&7&15&31&63&127&255&510&991&1024\\
\bottomrule
\end{tabular}
\end{center}
These dimensions are exact elimination results, not sampled action ranks.
The independent word certificate provides an upper bound of ten.

\begin{proof}[Proof of the length assertion in Theorem~\ref{thm:main}]
For any two generators, the number of words of length at most nine,
including the identity, is
\[
 1+2+\cdots+2^9=2^{10}-1=1023.
\]
Their span has dimension at most 1,023, whereas
$\dim_K\Mat_{32}(K)=1024$. Thus every generating pair requires length
at least ten. The 1,024-word certificate attains this bound over $\F$.
Its flattened basis matrix has determinant one, so it also remains a basis
over every field of characteristic two. Finally, an element of a binary
1,024-dimensional vector space is a unique XOR sum of its basis elements.
\end{proof}

The number ten measures the largest word length in a linear expansion.
It does not bound the number of words required in that expansion or the
running time of an implementation. The exact dimensions at depths eight
and nine also show that counting available words alone does not establish
independence.

\section{Invariant subspaces and quotients}
\begin{corollary}\label{cor:invariant}
Over every field $K$ of characteristic two, $A,B$ have no common invariant
subspace other than $0$ and $K^{32}$. Their common commutant is $K I$.
\end{corollary}
\begin{proof}
A subspace invariant under $A,B$ is invariant under all their products
and linear combinations, hence under every matrix. If it contains a
nonzero vector, arbitrary matrices send that vector to every vector of
$K^{32}$. It is therefore the whole space.
A matrix commuting with $A,B$ commutes with every matrix unit $E_{ij}$.
Commuting with all $E_{ii}$ makes it diagonal; commuting with all $E_{ij}$
makes its diagonal entries equal.
\end{proof}

\begin{corollary}\label{cor:quotient}
Let $Q:K^{32}\to K^r$ be linear. Suppose there are functions $f_A,f_B$,
not necessarily linear, such that
\[
 Q(Ax)=f_A(Qx),\qquad Q(Bx)=f_B(Qx)\quad\text{for all }x\in K^{32}.
\]
Then $\rank Q$ is either zero or 32. The same conclusion applies to the
linear part of an affine map with this property.
\end{corollary}
\begin{proof}
Evaluation at zero gives $f_A(0)=f_B(0)=0$. Hence $\ker Q$ is invariant
under both generators. Corollary~\ref{cor:invariant} makes this kernel
either the whole space or zero. For an affine map $Q(x)=Lx+c$, replace
$f_g$ by $y\mapsto f_g(y+c)-c$ and apply the same argument to $L$.
\end{proof}
This is a statement about a fixed linear or affine summary that supports
both operators on the whole input space. It leaves other types of
representation and restricted domains as separate questions.

\section{Verification and relation to earlier work}
The first certificate was recovered from an earlier exact computation by
the author. This publication supplies the explicit polynomial rank-one
certificate and the sharp length interpretation. The two certificates
use different checks. The generator forms matrices by input columns and
reduces from the highest pivot. The verifier derives the coordinate rows
directly from~\eqref{eq:shifts} and reduces from the lowest pivot.

The complete reconstruction experiment synthesizes all 1,024 elementary
maps $E_{ij}$ by both bases and compares every matrix entry. It also
reconstructs 64 additional binary matrices, checks both sides of all
inverse identities, and rejects five altered certificates. Equality on
the 32 input basis vectors proves equality of each complete linear map.
The release retains the exact certificates, arithmetic implementations,
and execution record. These are internal exact computations; a formal
proof-kernel check is a separate task.

The general rank-one mechanism is established matrix-algebra theory.
Radjabalipour, Rosenthal and Yahaghi discuss rank-one criteria over arbitrary
fields and division rings~\cite{radjabalipour}. The convention for word
length and the role of the identity are treated by Laffey, Markova and
\v{S}migoc~\cite{laffey}. Klep and \v{S}penko prove logarithmic word-length
bounds for generic generating families~\cite{klep}. Those general results
do not by themselves certify the fixed sparse pair in~\eqref{eq:operators}.
The contribution of this note is its explicit instance, optimal-length
certificate, rank-one identity, and reproducible full-map synthesis.

\appendix
\section{Inverse matrices for the cyclic bases}
Let $U,V$ be the matrices in~\eqref{eq:rankone}. Table~\ref{tab:inverse}
gives column $j$ of $X=U^{-1}$ and $Y=V^{-1}$ as eight-digit hexadecimal
words. The least significant bit is row zero. These 64 words, together
with~\eqref{eq:operators}, \eqref{eq:p} and~\eqref{eq:uv}, permit the
rank-one generation proof to be checked without the larger word list.
\begin{table}[H]
\centering\small
\begin{tabular}{rll@{\qquad}rll}
\toprule
$j$ & $X_j$ & $Y_j$ & $j$ & $X_j$ & $Y_j$ \\
\midrule
0 & \texttt{f3cbebea} & \texttt{20b2169f} & 16 & \texttt{910ddd88} & \texttt{a011bdee} \\
1 & \texttt{079451e0} & \texttt{b57e3ddc} & 17 & \texttt{c9b259b8} & \texttt{979b7737} \\
2 & \texttt{bb35cf67} & \texttt{3a9d3ea1} & 18 & \texttt{b983ccae} & \texttt{78ad91fe} \\
3 & \texttt{baee473e} & \texttt{9da2e048} & 19 & \texttt{5007ade7} & \texttt{9a438345} \\
4 & \texttt{c5fc4d3c} & \texttt{1a004498} & 20 & \texttt{99d1e586} & \texttt{68ef49fd} \\
5 & \texttt{92a2e5ca} & \texttt{f81a73cc} & 21 & \texttt{91e3c657} & \texttt{3f3745df} \\
6 & \texttt{acaac214} & \texttt{e56415f0} & 22 & \texttt{e40a2325} & \texttt{99410cb7} \\
7 & \texttt{d3e21a4a} & \texttt{c4edf9c1} & 23 & \texttt{4d3f6ee2} & \texttt{dfab6c2c} \\
8 & \texttt{b896618d} & \texttt{fcceda36} & 24 & \texttt{fa7584f7} & \texttt{e975ed5c} \\
9 & \texttt{ad581bc2} & \texttt{b36b0715} & 25 & \texttt{c73dd0ea} & \texttt{f50800fb} \\
10 & \texttt{c9b4c6cb} & \texttt{3c38bd92} & 26 & \texttt{52d75144} & \texttt{4a59035e} \\
11 & \texttt{dd23d427} & \texttt{a944f10a} & 27 & \texttt{e7761753} & \texttt{32f1698b} \\
12 & \texttt{3696d7a6} & \texttt{415bc077} & 28 & \texttt{4fb5683b} & \texttt{b599bca5} \\
13 & \texttt{413409c6} & \texttt{84811b0b} & 29 & \texttt{35ff5323} & \texttt{b2ca2295} \\
14 & \texttt{20ba532b} & \texttt{febbed6f} & 30 & \texttt{ce67fd01} & \texttt{893553fb} \\
15 & \texttt{5dfff284} & \texttt{a639bcc2} & 31 & \texttt{7a7a7285} & \texttt{bff0eed2} \\
\bottomrule
\end{tabular}
\caption{Exact inverse columns for the two cyclic bases.}
\label{tab:inverse}
\end{table}


\begin{thebibliography}{9}
\bibitem{rfc6234}
D. Eastlake 3rd and T. Hansen.
\textit{US Secure Hash Algorithms (SHA and SHA-based HMAC and HKDF)}.
RFC 6234, May 2011. \href{https://doi.org/10.17487/RFC6234}{doi:10.17487/RFC6234}.

\bibitem{radjabalipour}
M. Radjabalipour, P. Rosenthal and B. R. Yahaghi.
Burnside's theorem for matrix rings over division rings.
\textit{Linear Algebra and its Applications} \textbf{383} (2004), 29--44.
\href{https://doi.org/10.1016/j.laa.2003.08.011}{doi:10.1016/j.laa.2003.08.011}.

\bibitem{laffey}
T. Laffey, O. Markova and H. \v{S}migoc.
\textit{The effect of assuming the identity as a generator on the length of
the matrix algebra}. 2015.
\href{https://arxiv.org/abs/1501.05806}{arXiv:1501.05806}.

\bibitem{klep}
I. Klep and \v{S}. \v{S}penko.
\textit{Sweeping words and the length of a generic vector subspace of
$M_n(\mathbb F)$}. 2015.
\href{https://arxiv.org/abs/1506.03460}{arXiv:1506.03460}.
\end{thebibliography}
\end{document}
